\chapter{Fourier変換}
    \section{余弦,正弦変換}
    {%スコープ。\ge上書き回避
        \renewcommand{\ge}{g_\textrm{e}}
        \newcommand{\go}{g_\textrm{o}}
        \newcommand{\Ge}{G_\textrm{e}}
        \newcommand{\Go}{G_\textrm{o}}
        $f(t)$の余弦変換$C(\omega)$,正弦変換$S(\omega)$は次のように定義される。
        \begin{align*}
            C(\omega) &\coloneqq \sqrt{\frac{2}{\pi}}\integrate{0}{\infty}{f(x)\cos{\omega t}}{}{t}\\
            S(\omega) &\coloneqq \sqrt{\frac{2}{\pi}}\integrate{0}{\infty}{f(x)\sin{\omega t}}{}{t}
        \end{align*}
        $C,S$はそれぞれ\textcolor{red}{$\omega$の偶関数,奇関数}になっていることに注意。そして$t\geq0$において
        \begin{equation}
            f(t) \sim \sqrt{\frac{2}{\pi}}\integrate{0}{\infty}{C(\omega)\cos{\omega t}}{}{\omega} = \sqrt{\frac{2}{\pi}}\integrate{0}{\infty}{S(\omega)\sin{\omega t}}{}{\omega} \nonumber
        \end{equation}
        が成り立つ。注意すべきは、この逆変換は$t<0$に関しては何も教えてくれないということだ。変換時に$t<0$における情報を捨ててしまっているのだから当然である。
        以下、上の関係を示す。
        \newline\par
        \begin{proof}
            \quad\par
            $f(t)$の$t\geq0$の部分を取り出して$t<0$の部分に拡張することを考える。関数$\ge(t),\go(t)$を次のように定義する。
            \[\ge(t) \coloneqq f(|t|)\;\;(-\infty<t<\infty)\]
            \begin{align}
                \go(t) \coloneqq
                \begin{cases}
                    f(t) & (t>=0) \\
                    -f(-t) & (t<0)
                \end{cases}
                \nonumber
            \end{align}
            $\ge,\go$はそれぞれ$f$を偶関数,奇関数として拡張したものである。両者のフーリエ変換を$\Ge(\omega),\Go(\omega)$とすると
            \begin{align*}
                \Ge(\omega) &= \frac{1}{\sqrt{2\pi}}\integrate{-\infty}{\infty}{\ge(t)\NapierE^{-i\omega t}}{}{t} = \frac{1}{\sqrt{2\pi}}\integrate{-\infty}{\infty}{\ge(t)(\cos{\omega t}-i\sin{\omega t})}{}{t}\\
                &= \frac{1}{\sqrt{2\pi}}\integrate{-\infty}{\infty}{\ge(t)\cos{\omega t}}{}{t}\;\;\;(\because \ge(t)\sin{\omega t}\text{は奇関数なので対称積分の結果は0})\\
                &= \sqrt{\frac{2}{\pi}}\integrate{0}{\infty}{\ge(t)\cos{\omega t}}{}{t} = \sqrt{\frac{2}{\pi}}\integrate{0}{\infty}{f(t)\cos{\omega t}}{}{t} = C(\omega)
            \end{align*}
            となり余弦変換と一致する。同様にして
            \[\Go(\omega) = -iS(\omega)\]
            となることも確かめられる。
            \newline\par
            次に逆フーリエ変換で$\ge,\go$を復元してみる。
            \begin{align*}
                \ge(t) &\sim \frac{1}{\sqrt{2\pi}}\integrate{-\infty}{\infty}{\Ge(\omega)\NapierE^{i\omega t}}{}{\omega} = \frac{1}{\sqrt{2\pi}}\integrate{-\infty}{\infty}{C(\omega)\NapierE^{i\omega t}}{}{\omega} = \frac{1}{\sqrt{2\pi}}\integrate{-\infty}{\infty}{C(\omega)(\cos{\omega t}+i\sin{\omega t})}{}{\omega}\\
                &= \frac{1}{\sqrt{2\pi}}\integrate{-\infty}{\infty}{C(\omega)\cos{\omega t}}{}{\omega}\;\;(\because C(\omega)\sin{\omega t}\text{ は奇関数なので対称積分の結果は0})\\
                &= \sqrt{\frac{2}{\pi}}\integrate{0}{\infty}{C(\omega)\cos{\omega t}}{}{\omega}
            \end{align*}
            \begin{align*}
                \go(t) &\sim \frac{1}{\sqrt{2\pi}}\integrate{-\infty}{\infty}{\Go(\omega)\NapierE^{i\omega t}}{}{\omega} = \frac{1}{\sqrt{2\pi}}\integrate{-\infty}{\infty}{-iS(\omega)\NapierE^{i\omega t}}{}{\omega} = \frac{1}{\sqrt{2\pi}}\integrate{-\infty}{\infty}{-iS(\omega)(\cos{\omega t}+i\sin{\omega t})}{}{\omega}\\
                &= \frac{1}{\sqrt{2\pi}}\integrate{-\infty}{\infty}{S(\omega)\sin{\omega t}}{}{\omega}\;\;(\because S(\omega)\cos{\omega t}\text{ は奇関数なので対称積分の結果は0})\\
                &= \sqrt{\frac{2}{\pi}}\integrate{0}{\infty}{S(\omega)\sin{\omega t}}{}{\omega}
            \end{align*}
            $t\geq0$では$f,\ge,\go$は等しいから
            \[f(t) \sim \sqrt{\frac{2}{\pi}}\integrate{0}{\infty}{C(\omega)\cos{\omega t}}{}{\omega} = \sqrt{\frac{2}{\pi}}\integrate{0}{\infty}{S(\omega)\sin{\omega t}}{}{\omega}\]
        \end{proof}
        繰り返しになるが上式は$t\geq0$でしか成り立たない。$t<0$のことはわからないのである。クレイジーな変換だと思うかもしれないが、ちゃんと利点もある。実験データや音声,映像データは普通、記録開始を時刻0として扱う。この場合は$t<0$に関しては知る必要がない(どうでもいい)のだ。そういう状況下では複素数が登場する本家のフーリエ変換よりもsin,cosどっちか好きな方だけ使えば済むような正弦,余弦変換が役に立つことがある。
    }